\chapter{Case Study and Problem Statement}
\label{ch:casestudy}

\section*{Declaration of Contributions}
This chapter was prepared primarily by Nisarg Prajapati.
Nisarg Prajapati developed the XZ Utils case narrative, attack-path analysis,
and broader impact discussion. Adesh Palkar and Utkarsh Kumar assisted with
threat framing and technical interpretation.

\section{Case Study: XZ Utils Backdoor (2024)}

XZ Utils is a widely used compression utility in Linux ecosystems and includes
the \texttt{liblzma} library used by many software stacks.
In 2024, malicious code associated with releases 5.6.0 and 5.6.1 raised serious
concerns due to potential unauthorized access paths in environments linked with
\texttt{sshd}/OpenSSH \citep{nvd2024cve3094,redhat2024xz}.

\begin{table}[htbp]
\centering
\caption{Simplified timeline and implications of the XZ Utils incident.}
\label{tab:xz-timeline}
\begin{tabular}{p{0.22\textwidth} p{0.66\textwidth}}
\toprule
\textbf{Event} & \textbf{Observation} \\
\midrule
Upstream release activity & Malicious behavior was associated with release/build artifacts.
\\
Potential integration path & Systems using affected dependency paths could inherit risk transitively.
\\
Ecosystem response & Advisories and coordinated mitigation were issued to reduce spread.
\\
Security lesson & A trusted low-level dependency can become a high-impact attack surface.
\\
\bottomrule
\end{tabular}
\end{table}

\section{Attack Path and Security Insight}

The attack path can be interpreted as follows:

\begin{enumerate}
\item compromise is introduced at an upstream package/release stage;
\item downstream systems consume the package through normal dependency workflows;
\item malicious behavior is activated in operational contexts where trust is assumed.
\end{enumerate}

This pattern demonstrates a core insight: in supply chain attacks, \emph{implicit trust}
in dependency ecosystems can bypass traditional defensive boundaries.

\section{Problem Statement}

\textbf{The widespread blind trust in unverified open-source dependencies exposes organizations
to devastating software supply chain attacks.}

Technically, this problem appears because:

\begin{itemize}
\item modern software depends on deep transitive dependency chains;
\item package adoption is often based on popularity or convenience rather than verified provenance;
\item release artifact integrity and maintainer-risk checks are inconsistently enforced.
\end{itemize}

\section{Why This Matters}

A single compromised package can affect many products and organizations across sectors.
Attack impact propagates from upstream maintainers to downstream applications and then to end users.
National-level advisories (for example, CISA guidance) reinforce the need for rapid,
layered defense strategies and coordinated response \citep{cisa2024shieldsup}.

The solution direction is developed in \cref{ch:solutions}, where we compare current controls
and propose a composite mitigation strategy.
